% Part: incompleteness
% Chapter: incompleteness-provability
% Section: rosser-thm

\documentclass[../../../include/open-logic-section]{subfiles}

\begin{document}

\olfileid[pt]{inc}{inp}{ros}

\olsection{O teorema de Rosser}

Podemos modificar a prova de G\"odel para obter um resultado mais forte,
substituindo ``$\omega$-consistente'' simplesmente por ``consistente''? A
resposta é ``sim'', graças a um artifício descoberto por Rosser. O artifício
de Rosser consiste em usar um predicado de !!{derivability} ``modificado''
$\ORProv_T(y)$ em vez de $\OProv[\Th{T}](y)$.

\begin{thm}
\ollabel{thm:rosser}
Seja $\Th{T}$ qualquer teoria consistente e !!{axiomatizable} que
estenda $\Th{Q}$. Então $\Th{T}$ não é completa.
\end{thm}

\begin{proof}
Lembre-se de que $\OProv[\Th{T}](y)$ é definido como
$\lexists[x][\OPrf[\Th{T}](x, y)]$, onde $\OPrf[\Th{T}](x, y)$ representa
a relação decidível que vale sse $x$ é o número de G\"odel de
!!a{derivation} da !!{sentence} com número de G\"odel~$y$. A relação que
vale entre $x$ e~$y$ quando $x$~é o número de G\"odel de uma
\emph{refutação} da sentença com número de G\"odel~$y$ também é decidível.
Seja $\fn{not}(x)$ a função recursiva primitiva que faz o seguinte: se $x$
é o código de uma fórmula $!A$, $\fn{not}(x)$ é um código de $\lnot !A$.
Então $\Refut[\Th{T}](x, y)$ vale sse
$\Prf[\Th{T}](x, \fn{not}(y))$. Seja $\ORefut[\Th{T}](x, y)$ uma fórmula
que a represente. Assim, se $\Th{T} \Proves \lnot !A$ e $\delta$ é uma
!!{derivation} correspondente, $\Th{Q} \Proves
\ORefut[\Th{T}](\gn{\delta}, \gn{!A})$. Definimos $\ORProv[\Th{T}](y)$ como
\[
\lexists[x][(\OPrf[\Th{T}](x,y) \land \lforall[z][(z < x \lif \lnot
  \ORefut[\Th{T}](z,y))])].
\]
Em linhas gerais, $\ORProv[\Th{T}](y)$ diz ``há uma prova de $y$ em
$\Th{T}$ e não há refutação mais curta de~$y$''. Supondo que $\Th{T}$ seja
consistente, $\ORProv[\Th{T}](y)$ é verdadeira para os mesmos números que
$\OProv[\Th{T}](y)$; mas, do ponto de vista da \emph{demonstrabilidade}
em~$\Th{T}$ (e agora sabemos que há uma diferença entre verdade e
demonstrabilidade!), as duas têm propriedades distintas. Se $\Th{T}$ é
\emph{in}consistente, elas \emph{não} valem para os mesmos números!{}
($\ORProv[\Th{T}](y)$ é frequentemente lida como ``$y$ é demonstrável à
Rosser''. Como, conforme acabamos de discutir, a demonstrabilidade à Rosser
não é um tipo especial de demonstrabilidade---em teorias inconsistentes há
!!{sentence}s demonstráveis, mas não demonstráveis à Rosser---isso pode ser
confuso. Para evitar a confusão, poderíamos lê-la como ``$y$ é
shmovível''.)

Pelo lema do ponto fixo, existe uma fórmula $!R_\Th{T}$ tal que
\begin{equation}
  \Th{Q} \Proves !R_\Th{T} \liff \lnot \ORProv[\Th{T}](\gn{!R_\Th{T}}).
  \ollabel{RT}
\end{equation}
Em contraste com a prova de \olref[1in]{thm:first-incompleteness}, afirmamos
aqui que, se $\Th{T}$ é consistente, $\Th{T}$ não !!{derive}
$!R_\Th{T}$, e $\Th{T}$ tampouco !!{derive} $\lnot !R_\Th{T}$. (Em outras palavras,
não precisamos da hipótese de $\omega$-consistência.)

Primeiro, mostremos que $\Th{T} \Proves/ !R_{\Th{T}}$. Suponha que ela o
provasse, de modo que exista !!a{derivation} de~$!R_{\Th{T}}$ a partir
de~$T$; seja $n$ seu número de G\"odel. Então
$\Th{Q} \Proves \OPrf[\Th{T}](\num{n}, \gn{!R_{\Th{T}}})$, pois
$\OPrf[\Th{T}]$ representa $\Prf[\Th{T}]$ em~$\Th{Q}$. Além disso, para
cada $k < n$, $k$ não é o número de G\"odel de !!a{derivation} de
$\lnot !R_{\Th{T}}$, pois $\Th{T}$ é consistente. Logo, para cada $k < n$,
$\Th{Q} \Proves \lnot
\ORefut[\Th{T}](\num{k}, \gn{!R_{\Th{T}}})$. Por
\olref[req][min]{lem:less-nsucc},
$\Th{Q} \Proves \lforall[z][(z < \num{n} \lif \lnot \ORefut[\Th{T}](z,
  \gn{!R_{\Th{T}}}))]$. Portanto,
\[
\Th{Q} \Proves \lexists[x][(\OPrf[\Th{T}](x,\gn{!R_{\Th{T}}}) \land \lforall[z][(z
    < x \lif \lnot \ORefut[\Th{T}](z,\gn{!R_{\Th{T}}}))])],
\]
mas isso é precisamente $\ORProv[\Th{T}](\gn{!R_{\Th{T}}})$. Por
\olref{RT}, $\Th{Q} \Proves \lnot !R_{\Th{T}}$. Como $\Th{T}$ estende
$\Th{Q}$, também $\Th{T} \Proves \lnot !R_{\Th{T}}$. Supusemos que
$\Th{T} \Proves !R_{\Th{T}}$; portanto, $\Th{T}$ seria inconsistente,
contrariando a hipótese do teorema.

Agora mostremos que $\Th{T} \Proves/ \lnot !R_{\Th{T}}$. Suponha novamente
o contrário, e suponha que $n$ seja o número de G\"odel de
!!a{derivation} de~$\lnot !R_{\Th{T}}$. Então
$\Refut[\Th{T}](n, \Gn{!R_{\Th{T}}})$ vale e, como $\ORefut[\Th{T}]$
representa $\Refut[\Th{T}]$ em $\Th{Q}$, $\Th{Q} \Proves
\ORefut[\Th{T}](\num{n}, \gn{!R_{\Th{T}}})$. Mostraremos outra vez que
$\Th{T}$ seria então inconsistente, pois também !!{derive}ria
~$!R_{\Th{T}}$. Como
\begin{align*}
\Th{Q} & \Proves !R_{\Th{T}} \liff \lnot \ORProv[\Th{T}](\gn{!R_{\Th{T}}}), 
\intertext{e como $\Th{T}$ estende~$\Th{Q}$, basta mostrar que}
\Th{Q} & \Proves \lnot \ORProv[\Th{T}](\gn{!R_{\Th{T}}}).
\end{align*}
A !!{sentence} $\lnot
\ORProv[\Th{T}](\gn{!R_{\Th{T}}})$, O.s.\@,
\begin{align*}
  \lnot & \lexists[x][(\OPrf[\Th{T}](x,\gn{!R_{\Th{T}}}) \land
    \lforall[z][(z < x \lif \lnot \ORefut[\Th{T}](z,\gn{!R_{\Th{T}}}))])],
  \intertext{é logicamente equivalente a}
  & \lforall[x][(\OPrf[\Th{T}](x,\gn{!R_{\Th{T}}}) \lif
    \lexists[z][(z < x \land \ORefut[\Th{T}](z,\gn{!R_{\Th{T}}}))])].
\end{align*}
Argumentamos informalmente, usando a lógica e fatos sobre o que $\Th{Q}$
!!{derive}s. Suponha que $x$ seja arbitrário e que
$\OPrf[\Th{T}](x, \gn{!R_{\Th{T}}})$. Já sabemos que
$\Th{T} \Proves/ !R_{\Th{T}}$ e, portanto, para todo $k$,
$\Th{Q} \Proves \lnot \OPrf[\Th{T}](\num{k}, \gn{!R_{\Th{T}}})$. Logo,
para todo $k$, segue-se que $\eq/[x][\num{k}]$. Em particular, temos (a)
$\eq/[x][\num{n}]$. Também temos $\lnot(\eq[x][\num{0}]
\lor \eq[x][\num{1}] \lor \dots \lor \eq[x][\num{n-1}])$ e, por
\olref[req][min]{lem:less-nsucc}, (b) $\lnot(x < \num{n})$. Por
\olref[req][min]{lem:trichotomy}, $\num{n} < x$. Como $\Th{Q} \Proves
\ORefut[\Th{T}](\num{n}, \gn{!R_{\Th{T}}})$, temos $\num{n} < x \land
\ORefut[\Th{T}](\num{n}, \gn{!R_{\Th{T}}})$ e, daí, $\lexists[z][(z < x
\land \ORefut[\Th{T}](z,\gn{!R_{\Th{T}}}))]$. Como $x$ era arbitrário,
obtemos, como exigido,
\[
\lforall[x][(\OPrf[\Th{T}](x,\gn{!R_{\Th{T}}}) \lif
  \lexists[z][(z < x \land \ORefut[\Th{T}](z,\gn{!R_{\Th{T}}}))])].
\]
\end{proof}

\begin{prob}
Dois conjuntos $A$ e $B$ de números naturais são chamados
\emph{computavelmente inseparáveis} se não existe conjunto
!!{decidable}~$X$ tal que $A \subseteq X$ e
$B \subseteq \Complement{X}$ ($\Complement{X}$ é o complemento,
$\Nat \setminus X$, de~$X$). Seja $\Th{T}$ uma extensão consistente e
!!{axiomatizable} de $\Th{Q}$. Suponha que $A$ seja o conjunto dos números
de G\"odel das !!{sentence}s demonstráveis em~$\Th{T}$ e $B$ o conjunto
dos números de G\"odel das sentenças refutáveis em~$\Th{T}$. Prove que $A$
e~$B$ são computavelmente inseparáveis.
\end{prob}

\end{document}
